% ECA_16.tex
\documentclass[11pt,a4paper]{article}
\usepackage{amsmath,amssymb,amsthm}
\usepackage{geometry}
\geometry{margin=25mm}

\title{ECA\_16\\$\mathcal R^*$ の全部分集合に対する選択可能性}
\author{}
\date{}

\newtheorem{definition}{定義}
\newtheorem{proposition}{命題}
\newtheorem{remark}{注意}

\begin{document}
\maketitle

\section{目的}

ECA\_15では、可算無限公理候補族
\[
\mathcal R^*=\{r_n\mid n\in\mathbb N\}
\]
について、連続体濃度の議論へ進むためには、単に
\[
|\mathcal R^*|=\aleph_0
\]
であるだけでは不十分であり、その全部分集合をECAの選択空間 $I$ に接続する必要があることを整理した。

本稿では、このうち第2段階、

\[
\boxed{
\forall A\subseteq\mathcal R^*,
\quad
\mathcal A_{\mathrm{ZFC}}\cup A\in I
}
\]

を中心に検討する。

ここでは、この条件が現在のECAの定義から導かれるのか、それとも追加条件として導入する必要があるのかを区別する。

\section{既存の選択空間}

ECAでは、
\[
I\subseteq\mathcal P(\Ax)
\]
とする。

各
\[
i\in I
\]
はECAにおいて許容される一つの公理集合であり、
\[
i\subseteq\Ax
\]
を満たす。

また、
\[
B(i)\subseteq i\subseteq\Ax
\]
を維持する。

ここで注意すべきなのは、
\[
I\subseteq\mathcal P(\Ax)
\]
から
\[
I=\mathcal P(\Ax)
\]
は導かれないことである。

したがって、
\[
A\subseteq\Ax
\]
であることだけから
\[
A\in I
\]
を結論することはできない。

\section{宇宙別の選択可能性}

ECA\_14で導入した
\[
\operatorname{Ax}(U)\subseteq\Ax
\]
に対して、
\[
I(U)
=
\{i\in I\mid i\subseteq\operatorname{Ax}(U)\}
\]
とする。

したがって、
\[
I(U)\subseteq I.
\]

ここで、
\[
A\subseteq\operatorname{Ax}(U)
\]
であっても、
\[
A\in I(U)
\]
は自動的には成立しない。

実際、
\[
A\in I(U)
\iff
A\in I
\land
A\subseteq\operatorname{Ax}(U)
\]
である。

\section{$\mathcal R^*$ に対する選択可能性}

可算無限族
\[
\mathcal R^*\subseteq\operatorname{Ax}(U)
\]
を固定する。

ZFCの基準となる公理集合を
\[
\mathcal A_{\mathrm{ZFC}}\subseteq\operatorname{Ax}(U)
\]
とする。

このとき、任意の
\[
A\subseteq\mathcal R^*
\]
に対して
\[
\mathcal A_{\mathrm{ZFC}}\cup A
\]
は
\[
\operatorname{Ax}(U)
\]
の部分集合である。

しかし、
\[
\mathcal A_{\mathrm{ZFC}}\cup A
\subseteq\operatorname{Ax}(U)
\]
だけでは
\[
\mathcal A_{\mathrm{ZFC}}\cup A\in I
\]
を導けない。

したがって、次の条件を独立に考える。

\begin{definition}[$\mathcal R^*$ 選択閉包条件]
\[
\boxed{
\forall A\subseteq\mathcal R^*,
\quad
\mathcal A_{\mathrm{ZFC}}\cup A\in I
}
\]
を、$\mathcal R^*$ に対する選択閉包条件と呼ぶ。
\end{definition}

この条件は、
\[
\mathcal P(\mathcal R^*)
\]
全体をECAの選択空間へ写すための条件である。

\section{選択写像}

選択閉包条件が成立すると、
\[
\Gamma:
\mathcal P(\mathcal R^*)
\longrightarrow I
\]
を
\[
\boxed{
\Gamma(A)=\mathcal A_{\mathrm{ZFC}}\cup A
}
\]
と定義できる。

この写像は単射である。

実際、
\[
A\neq A'
\]
ならば、
\[
\mathcal A_{\mathrm{ZFC}}\cup A
\neq
\mathcal A_{\mathrm{ZFC}}\cup A'
\]
が必要になるためには、
\[
\mathcal R^*\cap\mathcal A_{\mathrm{ZFC}}=\varnothing
\]
という条件を置くのが安全である。

この条件の下では、$r\in A\triangle A'$ を取れば、
$r$ の所属によって両者が異なる。

したがって、
\[
\boxed{
\mathcal R^*\cap\mathcal A_{\mathrm{ZFC}}=\varnothing
\Longrightarrow
\Gamma\text{ は単射}
}
\]
である。

\section{重要な区別}

ここで、
\[
\mathcal P(\mathcal R^*)\hookrightarrow I
\]
が成立したとしても、
\[
\mathcal P(\mathcal R^*)\hookrightarrow\mathcal Z
\]
はまだ導かれない。

理由は、$I$ の要素は公理集合であり、
$\mathcal Z$ の要素はSolverであるためである。

したがって、
\[
\mathcal P(\mathcal R^*)
\hookrightarrow I
\]
から
\[
\mathcal P(\mathcal R^*)
\hookrightarrow\Solver_I
\]
へ進むには、各選択集合に対してSolverを構成する必要がある。

さらに、
\[
\mathcal P(\mathcal R^*)
\hookrightarrow\mathcal Z
\]
へ進むには、そのSolverがZFCを表現することを確認する必要がある。

\section{選択空間の全冪集合性との比較}

もし追加条件として
\[
I=\mathcal P(\Ax)
\]
を採用すれば、
\[
A\subseteq\Ax
\Longrightarrow
A\in I
\]
となる。

この場合、
\[
\mathcal A_{\mathrm{ZFC}}\cup A
\subseteq\Ax
\]
から
\[
\mathcal A_{\mathrm{ZFC}}\cup A\in I
\]
を直接得られる。

しかし、現在のECAは
\[
I\subseteq\mathcal P(\Ax)
\]
としており、
\[
I=\mathcal P(\Ax)
\]
を要求していない。

したがって、全冪集合性を新たに仮定することは、
現在のECAより強い条件となる。

\section{最小条件としての選択閉包}

連続体濃度の議論だけを目的とするなら、
\[
I=\mathcal P(\Ax)
\]
まで要求する必要はない。

必要なのは、少なくとも対象となる
\[
\mathcal P(\mathcal R^*)
\]
について、
\[
\mathcal A_{\mathrm{ZFC}}\cup A\in I
\]
がすべての
\[
A\subseteq\mathcal R^*
\]
に対して成立することである。

すなわち、
\[
\boxed{
\mathcal P(\mathcal R^*)
\text{ に対する局所的な選択閉包}
}
\]
で十分である。

これは、
\[
I=\mathcal P(\Ax)
\]
より弱い条件である。

\section{ECAから導出できるか}

現在の定義から、

\[
I\subseteq\mathcal P(\Ax)
\]

は与えられている。

しかし、この包含関係から、

\[
\forall A\subseteq\mathcal R^*,
\quad
\mathcal A_{\mathrm{ZFC}}\cup A\in I
\]

は導出できない。

したがって、

\[
\boxed{
\text{$\mathcal R^*$ 選択閉包条件は、現時点ではECAの定理ではない。}
}
\]

これはECA\_15で確認した第1段階の
\[
|\mathcal R^*|=\aleph_0
\]
とは独立した問題である。

\section{Solver構成との関係}

選択閉包条件が成立した場合、
各
\[
A\subseteq\mathcal R^*
\]
に対して
\[
i_A
=
\mathcal A_{\mathrm{ZFC}}\cup A
\in I
\]
を得る。

すると、
\[
\Solver_I
=
\{s=(\mathcal A_s,\mathcal R_s)\mid\mathcal A_s\in I\}
\]
において、
\[
\mathcal A_{s_A}=i_A
\]
を満たすSolver $s_A$ を構成することが次の課題となる。

したがって、構造は

\[
A
\longmapsto
i_A
\longmapsto
s_A
\longmapsto
\mathcal T(s_A)
\]

となる。

ここで各矢印は別々に検証する必要がある。

\section{濃度に関する帰結}

選択閉包条件と
\[
\mathcal R^*\cap\mathcal A_{\mathrm{ZFC}}=\varnothing
\]
が成立すれば、
\[
\mathcal P(\mathcal R^*)\hookrightarrow I
\]
を得る。

さらに
\[
|\mathcal R^*|=\aleph_0
\]
ならば、
\[
|\mathcal P(\mathcal R^*)|
=
2^{\aleph_0}.
\]

したがって、
\[
\boxed{
|I|\geq2^{\aleph_0}
}
\]
が得られる。

ただし、これは
\[
|\mathcal Z|\geq2^{\aleph_0}
\]
ではない。

$\mathcal Z$ の下界を得るためには、次に
\[
A\longmapsto s_A
\]
を構成し、さらに
\[
s_A\in\mathcal Z
\]
を確認する必要がある。

\section{結論}

第2段階の本質は、

\[
\boxed{
\mathcal P(\mathcal R^*)
\text{ の各要素をECAの選択空間 }I\text{ に保持できるか}
}
\]

という問題である。

現在の
\[
I\subseteq\mathcal P(\Ax)
\]
という定義だけでは、この性質は保証されない。

したがって、連続体濃度へ進むためには、

\[
\boxed{
\forall A\subseteq\mathcal R^*,
\quad
\mathcal A_{\mathrm{ZFC}}\cup A\in I
}
\]

を追加条件として採用するか、
あるいはECAの既存構造からこの条件を導出できることを示す必要がある。

また、
\[
I
\]
に対して全冪集合性
\[
I=\mathcal P(\Ax)
\]
を要求する方法もあるが、これは必要以上に強い。

したがって、現段階では
\[
\mathcal P(\mathcal R^*)\text{ に対する局所的な選択閉包}
\]
を最小条件として位置づける。

次の段階では、この選択された公理集合
\[
i_A=\mathcal A_{\mathrm{ZFC}}\cup A
\]
から実際に
\[
s_A\in\mathcal Z
\]
を構成できるか、すなわち第3段階である
\[
\boxed{
\forall A\subseteq\mathcal R^*,
\quad
\mathcal T(s_A)\cong_{\mathrm{Th}}\mathrm{ZFC}
}
\]
を検討する。

\end{document}
